\documentclass[11pt]{article}
\usepackage[margin=1in]{geometry}
\usepackage[T1]{fontenc}
\usepackage[utf8]{inputenc}
\usepackage{lmodern}
\usepackage{amsmath,amssymb,mathtools}
\usepackage{physics}
\usepackage{booktabs}
\usepackage{array}
\usepackage{enumitem}
\usepackage{xcolor}
\usepackage{hyperref}
\usepackage{microtype}
\usepackage{parskip}
\usepackage{titlesec}
\titleformat{\section}{\large\bfseries}{}{0em}{}
\titleformat{\subsection}{\normalsize\bfseries}{}{0em}{}
\hypersetup{colorlinks=true,linkcolor=blue,urlcolor=blue}
\newcommand{\kb}{k_\mathrm{B}}
\newcommand{\taueff}{\tau_\mathrm{eff}}
\begin{document}

\title{Module 2: From Radiation Interactions to Defect Generation in Silicon}
\author{}
\date{}
\maketitle

\section{Purpose of this module}
This module explains how the interaction history from the transport calculation is converted into a defect picture. In other words, Module 1 tells us how energy and momentum enter the silicon lattice; Module 2 tells us how that disturbance becomes vacancies, interstitials, and eventually electrically active defect states.

\section{The crystal picture}
Silicon in a good crystal can be pictured as a highly organized seating chart. Each atom has an expected seat and bonding arrangement. A radiation event that displaces an atom is like knocking one person out of their assigned seat. Immediately two things appear:
\begin{itemize}
    \item an empty seat, called a \textbf{vacancy},
    \item an extra misplaced atom, called an \textbf{interstitial}.
\end{itemize}
Together they form a Frenkel pair in the simplest picture.

\section{Primary knock-on atoms and cascade logic}
When an incoming particle transfers enough energy to a silicon atom, that atom becomes a \textbf{primary knock-on atom} (PKA). The PKA can then collide with neighboring atoms and create additional displacements. The simplest qualitative chain is:
\begin{equation}
\text{incident particle} \to \text{PKA} \to \text{secondary recoils} \to \text{point defects and clusters}.
\end{equation}

This matters because the final damage is often more than a single missing atom. It can become a small cascade of disorder.

\section{A first-order defect introduction model}
In the MATLAB model, the simplest assumption is that defect density scales with fluence:
\begin{equation}
N_D = k_D \Phi,
\end{equation}
where $N_D$ is defect concentration, $k_D$ is a defect introduction coefficient, and $\Phi$ is particle fluence.

\subsection{Why this equation makes sense}
This equation says that if each incoming particle has, on average, some fixed probability of producing stable defects, then doubling the number of particles should roughly double the number of defects. It is a first-order linear model.

A useful analogy is hail damaging a car roof. If each hailstone has some chance of making a dent, then more hailstones generally produce more dents. The exact dent pattern is complicated, but a linear first estimate is often reasonable at modest damage levels.

\subsection{What the coefficient means}
The coefficient $k_D$ absorbs many details:
\begin{itemize}
    \item particle type and energy,
    \item displacement efficiency,
    \item which defects survive rather than recombine immediately,
    \item how the chosen proxy is normalized.
\end{itemize}
So $k_D$ is not a fundamental constant in the deepest sense. It is a model parameter connecting irradiation conditions to stable defect concentration.

\section{Depth-dependent defect model}
If Geant4 gives a spatial damage proxy, then the model becomes
\begin{equation}
N_D(z) = k_D D_{\mathrm{proxy}}(z).
\end{equation}

This equation says the local defect concentration should be stronger where the transport calculation indicates stronger local damage.

\subsection{Interpretation}
The benefit of this form is that it respects nonuniform irradiation through the material. If the front side takes more damage than the back side, the defect profile should reflect that.

\section{Why defects become electrically important}
A structural defect matters electrically because it disturbs the periodic crystal potential. That disturbance can create allowed states inside the band gap. These states can:
\begin{itemize}
    \item trap carriers,
    \item assist recombination,
    \item assist generation current,
    \item scatter carriers and reduce mobility.
\end{itemize}

This is the bridge from lattice disorder to electronics. The crystal is no longer a perfect repeating environment, so carrier motion and occupancy are altered.

\section{Defect classes to keep in mind}
For this project, the exact spectroscopy of each defect is not required, but it helps to classify them qualitatively:
\begin{enumerate}[label=\arabic*.]
    \item \textbf{Point defects}: isolated vacancies and interstitials.
    \item \textbf{Complex defects}: vacancy-oxygen, vacancy-phosphorus, divacancies, and related combinations.
    \item \textbf{Clusters}: more disordered local regions created by stronger cascades.
\end{enumerate}

At the project level, these are grouped into an effective defect population $N_D$, but the physical reason that one effective parameter works is that many microscopic defects act through similar electrical roles: trapping, recombination, and generation.

\section{Derivation logic for the linear defect model}
A clean derivation uses rate balance. Suppose $d\Phi$ particles cross the material in a small increment. Let each increment create $\eta\, d\Phi$ stable defects per unit volume, where $\eta$ is an effective introduction efficiency. Then
\begin{equation}
\frac{dN_D}{d\Phi} = \eta.
\end{equation}
If $\eta$ is approximately constant, integration gives
\begin{equation}
N_D(\Phi) = N_D(0) + \eta \Phi.
\end{equation}
If the pre-irradiation defect background is neglected in the radiation-induced count, then
\begin{equation}
N_D \approx k_D \Phi.
\end{equation}
This derivation is useful because it shows the linear model is just the constant-rate solution of a simple balance law.

\section{When the linear model can fail}
The linear model is best viewed as the first term in a hierarchy. It may fail when:
\begin{itemize}
    \item annealing is important during irradiation,
    \item defect saturation occurs,
    \item cluster formation changes the efficiency at high dose,
    \item several defect species evolve with very different rates.
\end{itemize}

That is acceptable. For an interview-scale project, starting with the linear regime is both honest and useful.

\section{What to remember from this module}
\begin{itemize}
    \item A displacement event creates vacancies and interstitials.
    \item PKAs and cascades explain why damage can be richer than a single displaced atom.
    \item The first-order model $N_D = k_D \Phi$ is a constant-rate defect introduction law.
    \item Depth dependence can be added by mapping defect density to a Geant4-derived damage proxy.
    \item These defects become electrically important because they create trap, recombination, and generation centers.
\end{itemize}
\end{document}
